% \paragraph{Caption is the most important thing; image next. Paragraph is the least important---even our INLG paper said so! (CITATION). \cite{huang-etal-2023-summaries}}
% \paragraph{Caption is the most important thing; image next. Paragraph is the least important---even our INLG paper said so! (CITATION). \cite{huang-etal-2023-summaries}}

% \begin{figure}[t]
    \centering
    \includegraphics[width=\columnwidth]{figures/fig01/fig01_ablation_study-v5.png}
    %\caption{Result of ablation study on caption generation. Captions are the most important profile element, while images are more influential than paragraphs.
    %\kenneth{I like this design! (1) Enlarge font size. It's hard to read. (2) Can we pick a slightly darker color? These colors are fine, but a bit too light.}\sam{Updated with darker color and a bit larger font!}\kenneth{(1) What is the number on top of each bar? Maybe remove them. (2) Fonts need to be even bigger. (3) The label can be ``No Paragraph'', ``No Image''. ``No Caption'', (4) The bars can be slightly wider.}\sam{The number above is the average score for each setting. I think keeping them can make it more clear to know the actual difference? (Or create another table to show the full result of ablation study for clarity?)}\sam{P.S. I just make the chart flattened a bit, same setting as for EasyHard one}
    %}
    \vspace{-2pc}
    \caption{Ablation study on the impact of profile elements using GPT-4o. Results show a clear hierarchy of importance: caption > image > paragraph.}
    \vspace{-.5pc}
    \label{fig:ablation}
    %\vspace{-3mm}  % Reduce space
\end{figure}


%     \caption{Ablation study on the impact of profile elements. caption generation by removing one profile element at a time (caption, image, or figure-mentioning paragraph). Captions contribute the most, followed by images, then paragraphs.}
%\kenneth{TODO Human Evaluation here: https://openreview.net/forum?id=ghxUYlolwr\#discussion}


%\input{table/table_human rank}
\begin{figure}[t]
    \centering
    \includegraphics[width=\columnwidth]{figures/fig_humanRankDist/humanRankDist-v11.png}
    \vspace{-1.7pc}
    \caption{Distribution of human preference rankings (lower is better). A Friedman's test confirmed a statistically significant difference between the model configurations ($\chi^2$ = 28.48, $p < 0.001$).}
    \vspace{-1pc}
    \label{fig:human-rank}
    %\vspace{-3mm}  % Reduce space
\end{figure}

%\kenneth{(1) I'd use capitalized P for No-Profile. I know no-profile technically is one word, BUT I'd like to keep things consistent. (2) Can we try a bar chart (each model has one bar)? Like the one we always did for human evaluation.}\sam{corrected!}
%\kenneth{(1) Img resolution is a bit low. Use a better quality image. (2) The Rank 1, 2, 3, 4 legend/label is too close to the bar. Hard to see.}\sam{updated}

%\kenneth{I'd use All-Profile, No-Profile instead of All-profile, No-profile. Use upper case P, as they are basically names of conditions.}

%\kenneth{TODO: Which model do we use in human eval????}

%\sam{ADD-START}
%To supplement automatic metrics, which do not always correlate with quality, we conducted a human evaluation.
We recruited 10 US-based NLP researchers (PhD students with publication or review experience) for our human evaluation.
%\kenneth{Be more specific, are they PhD students? or?}\sam{corrected!} 
The study involved ranking captions for 50 figures randomly selected from the arXiv \textit{cs.CL} domain. 
For each figure, we provided the figure's image, the paper's title, and abstract as context. 
Participants then ranked four corresponding captions based on how well each one helped them understand the figure. The four captions were:
%\kenneth{What exactly is the QUESTION we asked the participants? Don't just say QUALITY.} \sam{corrected! Also included the full instruction in the Appendix as requested in the EMNLP checklist}
{\em (i)} the original \textit{Gold} caption, and
{\em (ii)} \textit{No-Profile}, 
{\em (iii)} \textit{1-Profile}, and 
{\em (iv)} \textit{All-Profile} settings generated by our best-performing model (\textbf{Gemini}, based on automatic metrics). 
See \autoref{app:human-eval} for details.
%The complete experimental setup, which follows prior work~\cite{hsu2025large}, is detailed in \autoref{app:human-eval}.
%with the complete experimental setup and participant instructions detailed in \autoref{app:human-eval}.
%The ranking results are presented in 

\paragraph{1-Profile was the most preferred condition.} Our human evaluation result shows a clear preference for \textit{1-Profile} setting, which achieved the best average rank of 2.27. The other models followed in order: \textit{No-Profile} (2.48) and \textit{All-Profile} (2.54). Interestingly, the author-written \textit{Gold} captions were ranked last overall with an average rank of 2.71. Such outcome is reflected in the preference distribution (\autoref{fig:human-rank}): The \textit{1-Profile} configuration received the most first-place votes (30.6\%) and the fewest last-place votes (17.8\%). On contrast, the \textit{Gold} captions were ranked last most frequently (31.2\%).
%\kenneth{TODOs: As we do not have a table now, write a few sentences here to report the avg rank for each model here.}\sam{done!}
Post-hoc Friedman-Nemenyi tests showed that 1-Profile and No-Profile captions were significantly preferred over gold ($p < 0.001$ and $p = 0.03$, respectively).

\paragraph{Trade-offs between human-perceived quality and similarity to gold captions.} 
While the All-Profile setting generated captions with higher reference similarity (\autoref{tab:llm_comparison}), human judges significantly preferred captions from 1-Profile ($p = 0.006$). 
This suggests that optimizing for similarity may also reproducing flaws from the inconsistently quality's caption from arXiv reference~\cite{huang-etal-2023-summaries}, reducing the perceived quality.

%This may stem from the variable quality of the reference arXiv captions themselves~\cite{huang-etal-2023-summaries}. 
%Consequently, when a model too closely mimics the gold standard, it may inherit stylistic flaws, causing perceived quality to decline as similarity increases excessively.

%\sam{ADD-END}
%------------ dead kitten ----------
\begin{comment}


    %While the All-Profile setting generated captions with higher reference similarity, with higher similarity to the references (\autoref{tab:llm_comparison}), human judges significantly preferred captions from 1-Profile ($p = 0.006$). 

    % This likely reflects the inconsistent quality of arXiv captions~\cite{huang-etal-2023-summaries}.
Models that closely mimic these references risk reproducing their flaws, reducing perceived quality as similarity increases.
\end{comment}



